% =====================================================================
%   Sensor-Only Wear Prediction on MATWI — Technical Report
%   Project PINN (Stage 2 sub-track)
%
%   Compile with:  pdflatex sensor_only_report.tex
% =====================================================================

\documentclass[11pt,a4paper]{article}

\usepackage[utf8]{inputenc}
\usepackage[T1]{fontenc}
\usepackage[margin=2.4cm]{geometry}
\usepackage{amsmath,amssymb}
\usepackage{booktabs}
\usepackage{array}
\usepackage{xcolor}
\usepackage{listings}
\usepackage{enumitem}
\usepackage{hyperref}
\hypersetup{
    colorlinks=true,
    linkcolor=blue!60!black,
    urlcolor=blue!60!black,
    citecolor=blue!60!black,
}

% --- Listings style for file / code references --------------------------
\definecolor{codebg}{rgb}{0.96,0.96,0.96}
\lstset{
    basicstyle=\ttfamily\small,
    backgroundcolor=\color{codebg},
    breaklines=true,
    showstringspaces=false,
    columns=fullflexible,
    frame=single,
    framerule=0pt,
    xleftmargin=4pt,
    framexleftmargin=4pt,
}

% --- Convenience commands -----------------------------------------------
\newcommand{\code}[1]{\texttt{#1}}
\newcommand{\um}{\,\textmu m}

\title{Sensor-Only Wear Prediction on MATWI \\[2pt]
       \large Engineered features, per-set deltas, and gradient boosting}
\author{Project PINN --- Stage 2 sub-track}
\date{}

\begin{document}
\maketitle

\begin{abstract}
\noindent
This report documents the sensor-only sub-track of the MATWI wear-prediction
project. The starting point was a Ridge regression on 40 hand-engineered
sensor features (\code{sensor\_check.py}), which produced a test
MAE of 56.0\um{} --- \emph{worse} than predicting the training mean (40.9\um).
A second-generation pipeline (\code{sensor\_only\_v2.py}) was built that
combines (i)~a richer dimensionless feature set, (ii)~per-set delta features
that subtract a learned ``unworn baseline,'' (iii)~LightGBM in addition to
Ridge, and (iv)~leave-one-set-out cross-validation. The best resulting
method, \code{lgbm\_v2cut\_absolute}, reaches \textbf{29.9\um{} test MAE}
($-11\um$ vs.\ predict-mean, $-26\um$ vs.\ the original baseline) with
\textbf{54.8\um{} LOSO-pooled MAE}. All LightGBM variants cluster in the
$29$--$33\um$ test band, indicating a genuine ceiling on the engineered-feature
sensor-only setup. The dominant remaining error source is two anomalous
training sets (5 and 11) whose LOSO folds blow up to 93--169\um{} regardless
of model or features, and which are \emph{not} explained by cutting-parameter
differences. The result is a defensible sensor-only baseline; further sensor
work shows diminishing returns relative to the project's main contribution,
the Stage 3 physics-informed loss.
\end{abstract}

% =====================================================================
\section{Motivation and scope}
% =====================================================================

The MATWI dataset (De Pauw et al., KU Leuven, 2023) provides paired
images and multi-channel sensor recordings for 17 milling tools run to
failure. The project's main pipeline performs \emph{vision-only} and
\emph{vision+sensor fusion} wear prediction. Before invoking fusion, a
question worth asking is: \textbf{can the sensor channels alone predict
wear at all?}

A first attempt (\code{sensor\_check.py}, the predecessor of this
sub-track) answered: not really. A Ridge regression on 40 hand-engineered
sensor features achieved 56.0\um{} test MAE --- a model that performs
\emph{worse} than the trivial baseline of always predicting the training
mean (40.9\um). The training MAE of 21.7\um{} confirmed massive
overfitting; the model was memorising set identity, not learning wear.

This sub-track set out to build a meaningful sensor-only baseline. The
goal was not to match vision-only performance (19\um{} test MAE), which
is unrealistic on this dataset, but to (a)~beat predict-mean by a credible
margin, (b)~understand \emph{why} the v1 attempt failed, and (c)~establish
a clean reference for future physics-informed and fusion experiments.

% =====================================================================
\section{Folder contents}
% =====================================================================

All files described below live in this folder (\code{sensor-only/}). The
shared multimodal dataset class is in the sibling \code{vision-sensor/}
folder and is imported via a path-search helper.

\begin{description}[style=nextline,labelindent=0pt,leftmargin=2em]

\item[\code{sensor\_check.py}]
    The first-generation sensor-only baseline. Reuses the
    \code{MATWIMultimodalDataset} class from \code{vision-sensor/} to
    extract its 40 hand-engineered features per sample (5 channels
    $\times$ 8 features: mean, std, RMS, peak-to-peak, kurtosis,
    dominant frequency, low-band energy, mid-band energy). Fits a Ridge
    regression with an L2 standardiser; reports MAE on train, val, test
    broken down by wear type, and dumps per-feature
    $|$coefficient$|$ rankings. Used to diagnose the original failure
    mode (it overfits to absolute, gain-dependent spectral energies).

\item[\code{sensor\_check.sh}]
    SLURM submission script for \code{sensor\_check.py}, configured for
    Snellius (gpu\_a100, 30~minutes, 4~CPUs, 16~GB). The CPU budget is
    the bottleneck because feature extraction is dominated by CSV I/O,
    not GPU compute.

\item[\code{sensor\_only\_v2.py}]
    The second-generation pipeline and the main contribution of this
    sub-track. Re-implements the feature extractor with $\sim$85 dimensionless
    features (relative spectral bands, robust spread statistics,
    within-pass windowed RMS trend, cross-channel force resultant);
    adds an optional per-set delta transform; supports both Ridge and
    LightGBM models; loads cutting parameters (\code{Vc, fz, Vf, Ae,
    Ap, z, material}) from \code{sets.csv} as additional features for
    LightGBM; and runs all comparisons under both the official
    train/val/test split and leave-one-set-out cross-validation on
    training sets. Outputs are cached so re-runs take $\sim$30 seconds
    instead of 5 minutes.

\item[\code{sensor\_only\_report.tex}]
    This document.

\end{description}

Results from running \code{sensor\_only\_v2.py} are written by default to
\code{runs/sensor\_only\_v2/} at the project root:

\begin{lstlisting}
runs/sensor_only_v2/
  features_cache_1_13.npz     extracted v1+v2 features (reused on reruns)
  results.json                full per-method, per-split metrics
  results.csv                 flat table (sorted by test MAE)
  loso_results.csv            per-fold leave-one-set-out breakdown
  feature_importance.csv      top 20 features of the best method
\end{lstlisting}

% =====================================================================
\section{Methodology}
% =====================================================================

% ----------------------------------------------------------------
\subsection{Diagnosis of the v1 failure}
% ----------------------------------------------------------------

The v1 extractor's 40 features included \emph{absolute} spectral band
energies (\code{Fx\_low\_energy = np.sum(fft\_power[mask])}). These
quantities scale linearly with sensor gain, signal amplitude, and
recording length --- all of which differ per set in MATWI. The v1
Ridge model's top-5 features by $|\text{coef}|$ were exactly these
band energies, meaning the model was using them as set-identity
proxies. When the train$\to$test distribution shift kicked in (train
is CK45 only; val/test include RVS\,304), those proxies became
anti-correlated with wear and the model collapsed.

Three concrete problems were identified in the v1 extractor:

\begin{enumerate}
    \item \textbf{Absolute spectral energies} (gain- and length-dependent).
    \item \textbf{Whole-pass global statistics} (\code{p2p}, \code{kurtosis})
          over $\sim$99{,}000 samples, dominated by single outliers; whole-pass
          dominant frequency, which is essentially meaningless on a broad
          spectrum.
    \item \textbf{No cross-channel physical features}, although tool wear
          physically raises the cutting force resultant $F_\text{res} =
          \sqrt{F_x^2 + F_y^2 + F_z^2}$ in a coordinate-frame-independent
          way.
\end{enumerate}

% ----------------------------------------------------------------
\subsection{Tier 1 --- v2 feature engineering}
% ----------------------------------------------------------------

The new feature extractor produces 85 features per sample:

\paragraph{Per channel (15 features $\times$ 5 channels = 75).}
For each of \code{acc, acoustic, Fx, Fy, Fz}:
\begin{description}[style=nextline,labelindent=0pt,leftmargin=2em]
    \item[Basic statistics (3)] mean, std, RMS.
    \item[Robust spread (2)] interquartile range $p_{75}-p_{25}$ and
        $p_{95}-p_{5}$. Replaces the v1 \code{p2p}, which over 99k
        samples reduces to a single outlier.
    \item[Robust shape (1)] fraction of samples beyond $3\sigma$ from
        the mean. Replaces the v1 \code{kurtosis}, which is similarly
        outlier-dominated.
    \item[Windowed RMS trend (4)] split the pass into $N{=}10$
        contiguous windows; compute RMS per window; fit a line to RMS
        vs.\ normalised window index. Report (mean, std, slope,
        intercept). The \emph{slope} is the within-pass temporal
        evolution that whole-pass globals destroy.
    \item[Spectral (5)] three \emph{relative} band energies (each
        band's power divided by total power), spectral entropy
        (Shannon, on the normalised power spectrum), and spectral
        centroid in Hz. All five are dimensionless or in physical
        units that do not depend on sensor gain or recording length.
\end{description}

\paragraph{Cross-channel (10 features).}
\begin{itemize}
    \item Force resultant $F_\text{res}$: mean, std, windowed
          (mean\_rms, std\_rms, slope\_rms, intercept\_rms). Six features.
    \item Mean of $|F_x|/|F_y|$ ratio.
    \item Per-axis coefficient of variation $\sigma/|\mu|$ for $F_x,
          F_y, F_z$.
\end{itemize}

The hypothesis was that these features would generalise across setups
where the v1 features could not. Empirically, this turned out to be
true only for tree-based models (see Results).

% ----------------------------------------------------------------
\subsection{Tier 2 --- per-set delta features}
% ----------------------------------------------------------------

The MATWI dataset suffers from setup-to-setup signal-amplitude drift
that survives even careful feature engineering. To remove this for
linear models, the v2 pipeline computes a \emph{per-set baseline
feature vector} as the mean over the $K{=}3$ lowest-\code{SensorID}
passes for each set, then subtracts that baseline from every sample
in the set:
\begin{equation*}
    \tilde{x}_i \;=\; x_i \;-\;
        \frac{1}{K}\sum_{j \in \text{first-K}(s_i)}\!\!\!\! x_j .
\end{equation*}
The key observation: \code{SensorID} is the recording-order index
\emph{not} the wear label, so using the early-passes-as-reference is
legitimate even on val/test sets --- there is no label leakage.
Each set is also entirely contained in one official split (e.g.,
all of set~3 is in val), so the baseline computation does not cross
the split boundary.

Per-set deltas dramatically improved Ridge (test MAE 56.0$\to$34.8 for
the v1 features) but slightly hurt LightGBM, because tree splits do
not benefit from mean-centering --- subtracting a baseline destroys
absolute-scale information that splits could otherwise use.

% ----------------------------------------------------------------
\subsection{Tier 3 --- LightGBM}
% ----------------------------------------------------------------

LightGBM (Microsoft's gradient-boosted decision-tree implementation) is
the natural model class for tabular problems of this size. Compared to
Ridge, it captures nonlinear feature interactions and is invariant to
feature scaling (tree splits are threshold comparisons). Compared to
neural networks, it is far less prone to overfitting on $\sim$647
training samples and provides interpretable per-feature gain
importances.

Hyperparameters (small-data conservative):
\begin{lstlisting}
objective         = "regression_l1"   # MAE objective
learning_rate     = 0.03
num_leaves        = 31
max_depth         = 5
min_data_in_leaf  = 10
feature_fraction  = 0.8               # per-tree feature subsample
bagging_fraction  = 0.8               # per-tree row subsample
lambda_l2         = 1.0               # L2 on leaf values
num_boost_round   = 3000
early_stopping    = 150               # patience on val MAE
\end{lstlisting}

Missing values (e.g., Set~1's unknown cutting parameters) are passed
through natively --- LightGBM learns at each split whether NaN values
should be routed left or right.

The original v1 \code{sensor\_check.py} additionally provided a
\textbf{Ridge sanity check} that this sub-track did not replace.
Treating that script as the v1 baseline keeps the comparison honest:
both v1 Ridge and v2 Ridge are evaluated in \code{sensor\_only\_v2.py}
under identical splits and evaluation harnesses.

% ----------------------------------------------------------------
\subsection{Cutting-parameter features}
% ----------------------------------------------------------------

The eight cutting parameters from \code{sets.csv} ($V_c$, $n$, $f_z$,
$V_f$, $A_e$, $A_p$, $z$, material as a binary RVS flag) are appended
to the v2 features for the LightGBM \code{*cut*} methods. The csv has
a few messy cells: Set~1 has all-unknown parameters (encoded as
``?'') and Set~6 has variable feed rate ``185/148'' --- both handled
by the parser (\code{?}~$\to$~NaN, ``a/b''~$\to$~mean of parts).

The hypothesis was that adding cutting parameters would let LightGBM
\emph{route} predictions by cutting regime, especially for the
LOSO-outlier sets 5 and 11. As reported below, this hypothesis
failed empirically.

% ----------------------------------------------------------------
\subsection{Evaluation: official split + leave-one-set-out CV}
% ----------------------------------------------------------------

The standard evaluation is on the official train (sets
\{1,2,5,7,8,10,11\}) / val (\{3,6,12\}) / test (\{4,9,13\}) split,
where wear MAE is also broken down by wear type (flank, adhesion,
flank+adhesion). However, the official val set contains set~12
(RVS\,304, never seen in training) and is therefore an unreliable
guide to in-distribution generalisation.

\textbf{Leave-one-set-out (LOSO) CV} on the seven training sets
provides a fairer signal. For each training set $s$, train on the
other six sets and evaluate on $s$. Predictions are pooled across
folds and the pooled MAE is reported. For delta-based methods, the
per-set baseline is recomputed inside each fold (using the held-out
set's own \code{SensorID} order, so there is no label leakage).

The discrepancy between LOSO MAE and test MAE is itself diagnostic:
similar values mean the model generalises across setups; a much
worse LOSO indicates dependence on specific training sets.

% =====================================================================
\section{Results}
% =====================================================================

% ----------------------------------------------------------------
\subsection{Comparison matrix}
% ----------------------------------------------------------------

Table~\ref{tab:results} summarises all methods on the official split
plus LOSO-pooled CV on training sets.

\begin{table}[h]
\centering
\caption{Test, val and LOSO-pooled MAE in \textmu m. The full method
matrix from \code{results.csv}.}
\label{tab:results}
\small
\begin{tabular}{lrrrr}
\toprule
Method                          &  val  &  test &  LOSO  &  train \\
\midrule
predict\_mean\_train            & 57.7  & 40.9  &   ---  &  ---   \\
\midrule
ridge\_v1\_absolute (baseline)  & 85.8  & 56.0  &  78.7  &  21.7  \\
ridge\_v1\_delta                & 78.4  & 34.8  & 112.7  &  23.0  \\
ridge\_v2\_absolute             & 84.2  & 59.2  & 106.1  &  19.9  \\
ridge\_v2\_delta                & 88.6  & 50.1  &  69.0  &  18.7  \\
\midrule
lgbm\_v2\_absolute              & 49.9  & 31.9  &  55.0  &  49.3  \\
lgbm\_v2\_delta                 & 49.7  & 35.0  &  62.7  &  54.5  \\
lgbm\_v2cut\_absolute           & 50.2  & \textbf{29.9} & \textbf{54.8} &  44.8 \\
lgbm\_v1v2\_absolute            & 49.8  & 31.2  &  58.3  &  49.1  \\
lgbm\_v1v2cut\_absolute         & 49.9  & 32.6  &  60.9  &  51.3  \\
\bottomrule
\end{tabular}
\end{table}

\noindent
Five observations:

\begin{enumerate}
    \item \textbf{Best method: \code{lgbm\_v2cut\_absolute} at 29.9\um{}
          test, 54.8\um{} LOSO.} A clean $-11\um$ vs.\ predict-mean
          and $-26\um$ vs.\ the original Ridge baseline.
    \item \textbf{LightGBM variants cluster tightly (29.9--32.6\um{}
          test).} Cutting parameters, kitchen-sink feature unions,
          and patience tweaks all move the needle by less than
          $3\um$. The model is at its ceiling.
    \item \textbf{Per-set delta helps Ridge but hurts LightGBM.}
          Ridge needs the linear-model fix; tree splits handle scale
          natively and lose information when it is removed.
    \item \textbf{The v2 features helped only when combined with
          LightGBM.} Under Ridge, v2 was no better than v1 (in
          fact slightly worse) --- 85 features on 647 samples
          overfits Ridge. Under LightGBM, the new robust-spread,
          spectral, and cross-channel features are heavily used (see
          \S\ref{sec:fimp}).
    \item \textbf{Val is harder than test.} Set~12 (RVS\,304) makes
          val a true out-of-distribution split. Best-LightGBM val
          MAE (49.9\um) is barely below predict-mean on val (57.7\um),
          while test MAE drops by $\sim$11\um. The official val set
          is therefore unreliable for selecting between models;
          LOSO is the better signal.
\end{enumerate}

% ----------------------------------------------------------------
\subsection{Feature importance (best method)}
\label{sec:fimp}
% ----------------------------------------------------------------

Table~\ref{tab:fimp} lists the top 20 features of
\code{lgbm\_v2cut\_absolute} by LightGBM ``gain'' importance.

\begin{table}[h]
\centering
\caption{Top 20 features by LightGBM gain importance.
Eight of the top 20 are v2-specific (marked \textsuperscript{$\star$});
all ten cutting parameters were available but only one (\code{material\_RVS},
not shown) appeared in the top 30.}
\label{tab:fimp}
\small
\begin{tabular}{rlr}
\toprule
Rank & Feature & $|$gain$|$ \\
\midrule
 1 & \code{Fx\_p95mp05}\textsuperscript{$\star$}            & 1047 \\
 2 & \code{Fy\_std}                                          &  863 \\
 3 & \code{Fx\_std}                                          &  725 \\
 4 & \code{Fy\_p95mp05}\textsuperscript{$\star$}            &  242 \\
 5 & \code{Fy\_spec\_entropy}\textsuperscript{$\star$}      &  160 \\
 6 & \code{Fy\_mean}                                         &  155 \\
 7 & \code{Fx\_mean}                                         &  121 \\
 8 & \code{Fx\_cv}\textsuperscript{$\star$}                 &  106 \\
 9 & \code{Fx\_spec\_entropy}\textsuperscript{$\star$}      &   86 \\
10 & \code{Fy\_rel\_e\_b2}\textsuperscript{$\star$}         &   70 \\
11 & \code{Fy\_cv}\textsuperscript{$\star$}                 &   66 \\
12 & \code{Fx\_win\_mean\_rms}\textsuperscript{$\star$}     &   57 \\
13 & \code{acc\_frac3sig}\textsuperscript{$\star$}          &   57 \\
14 & \code{Fres\_win\_std\_rms}\textsuperscript{$\star$}    &   56 \\
15 & \code{Fy\_win\_mean\_rms}\textsuperscript{$\star$}     &   49 \\
16 & \code{Fy\_win\_intercept\_rms}\textsuperscript{$\star$}&   46 \\
17 & \code{Fz\_mean}                                         &   36 \\
18 & \code{acoustic\_frac3sig}\textsuperscript{$\star$}     &   34 \\
19 & \code{acc\_win\_std\_rms}\textsuperscript{$\star$}     &   32 \\
20 & \code{FxFy\_ratio\_mean}\textsuperscript{$\star$}      &   30 \\
\bottomrule
\end{tabular}
\end{table}

\noindent
The picture is consistent and physical: force channels dominate (17 of
the top 20), the force resultant cross-channel feature is used, and
the v2-specific robust spread features ($p_{95}{-}p_{5}$) hold the top
positions. Cutting parameters were available but not heavily used,
explaining why they added almost no test-MAE improvement.

% ----------------------------------------------------------------
\subsection{Per-fold LOSO breakdown}
\label{sec:loso}
% ----------------------------------------------------------------

Pooled LOSO MAE hides important structure. Table~\ref{tab:loso} shows
per-fold MAE for the two best methods.

\begin{table}[h]
\centering
\caption{Leave-one-set-out per-fold test MAE in \textmu m.
Sets 5 and 11 are systematic outliers --- adding cutting parameters
barely changes them.}
\label{tab:loso}
\small
\begin{tabular}{crrr}
\toprule
Heldout set & lgbm\_v2\_absolute & lgbm\_v2cut\_absolute & $\Delta$ \\
\midrule
 1 &  27.9 & 26.2 & $-1.7$ \\
 2 &  56.0 & 57.6 & $+1.6$ \\
 5 & 168.9 & \textbf{168.5} & $-0.4$ \\
 7 &  25.3 & 27.1 & $+1.8$ \\
 8 &  25.4 & 26.7 & $+1.3$ \\
10 &  28.9 & 25.9 & $-3.0$ \\
11 &  94.4 & \textbf{93.4} & $-1.0$ \\
\bottomrule
\end{tabular}
\end{table}

\noindent
\textbf{Five of seven folds give 25--57\um{} --- a perfectly respectable
sensor-only result.} The remaining two (5 and 11) blow up by a factor
of $3$--$6\times$, dragging the pooled LOSO from $\sim$28\um{} (if
those two sets were excluded) up to 55\um{}. Crucially, adding cutting
parameters does \emph{not} fix sets 5 and 11; the largest improvement
on either is $-1.7\um$. Sets 5 ($V_c{=}174, f_z{=}0.04, V_f{=}148$) and
11 ($V_c{=}174, f_z{=}0.043, V_f{=}159$) sit \emph{inside} the
parameter envelope of the other training sets, so they are not
extrapolating in regime space. The anomaly must therefore be in the
sensor signatures or wear-progression curves themselves --- likely
tool variation, lubrication, or sensor calibration drift on those
recordings --- and is not recoverable through feature engineering on
the existing channels.

% =====================================================================
\section{Findings}
% =====================================================================

\begin{enumerate}
    \item The original 56\um{} Ridge failure was caused by the
          extractor relying on \emph{absolute} spectral band energies
          that doubled as set-identity proxies. Replacing them with
          dimensionless spectral shape features (relative band energy,
          entropy, centroid) fixed this in principle, but the gain only
          materialised once the model class was changed too.

    \item \textbf{The single biggest intervention for sensor-only is
          switching to LightGBM.} Ridge with v2 features peaks at
          50\um{} test; LightGBM with the same features peaks at
          32\um{}. Tree splits both (a)~capture nonlinear interactions
          between features that Ridge cannot, and (b)~handle the
          residual setup-to-setup scale drift that Ridge stumbles on.

    \item \textbf{Per-set delta features are a Ridge-specific patch,
          not a universal one.} They cut Ridge test MAE on v1 features
          by $\sim$21\um{} but cost LightGBM $\sim$5\um. This is a useful
          lesson for any future modelling: linear-model fixes do not
          always transfer to tree models.

    \item \textbf{Cutting parameters helped slightly but did not
          explain the LOSO outliers.} The hypothesis that sets 5 and
          11 were anomalous because of their cutting regime was
          \emph{falsified}: the model can see the parameters and
          still fails on those folds. The actual cause is at the
          sensor or wear-progression level, not the cutting physics.

    \item \textbf{There is a real ceiling.} All LightGBM variants are
          within 3\um{} of each other on test and within 6\um{} on
          LOSO. The kitchen-sink combination (v1$\cup$v2$\cup$cutting)
          is in fact the worst LightGBM variant, indicating diminishing
          returns from feature accumulation. Further engineering on
          the same sensor channels is unlikely to move the needle.
\end{enumerate}

% =====================================================================
\section{Limitations and the ceiling}
% =====================================================================

The remaining error budget breaks down approximately as follows
(qualitative, not from a formal decomposition):

\begin{itemize}
    \item \textbf{Two anomalous training sets (5, 11)} account for the
          gap between pooled LOSO ($\sim$55\um) and the typical-fold
          LOSO ($\sim$28\um). These would require either dropping
          (yielding a more optimistic but smaller-data estimate) or a
          fundamentally different model class --- per-sample rather
          than per-pass features, for instance.

    \item \textbf{CK45 $\to$ RVS\,304 distribution shift} caps val and
          part of test performance. Set~12 (val) and set~13 (test) are
          RVS\,304 but no RVS\,304 appears in train under the official
          \code{set\_range="1-13"} setting. Running with
          \code{--set-range 1-17} (adding sets 14--17 to training)
          would address this but breaks comparability with the
          published baselines.

    \item \textbf{Sample size (647 training samples)} caps model
          complexity. Approaches like 1D-CNNs over raw downsampled
          sequences, or transformer encoders over windowed segments,
          would add learned representations on top of the engineered
          features, but each additional parameter risks overfitting
          on a corpus this small.

    \item \textbf{Vision-only is structurally better-suited to this
          task.} The vision-only baseline reaches 19\um{} test MAE on
          the same split because tool wear is fundamentally a visual
          phenomenon (a worn flank looks different from a fresh one),
          while the sensor signals reflect wear only indirectly through
          forces and vibrations that depend on many other factors.
          \textbf{Sensor-only is not expected to match vision-only on
          this dataset}, and reaching 29.9\um{} ($1.6\times$ vision-only)
          on a per-feature-engineering pipeline is a respectable
          result.
\end{itemize}

% =====================================================================
\section{How to reproduce}
% =====================================================================

\begin{lstlisting}[language=bash]
# From the project root, after the dataset is downloaded into
# dataset/matwi/ via dataset/download_data.py.

# Optional but recommended for the full method matrix:
pip install lightgbm

# Run all methods, save outputs under runs/sensor_only_v2/
python sensor-only/sensor_only_v2.py \
    --data-dir   dataset/matwi \
    --labels-csv dataset/matwi/labels.csv \
    --output-dir runs/sensor_only_v2

# First run extracts features (~5 min, dominated by CSV I/O).
# Subsequent runs reuse the cache and finish in ~30 s.

# To exclude cutting-parameter methods:
#   --no-cutting-params

# To use all sensor samples (drops the require-image filter that
# matches the existing 647-sample baseline):
#   --all-sensor-samples

# To include RVS 304 (sets 14-17) in training:
#   --set-range 1-17
\end{lstlisting}

% =====================================================================
\section{Conclusion}
% =====================================================================

The sensor-only sub-track produced a credible baseline
(\textbf{29.9\um{} test, 54.8\um{} LOSO}) and a clear diagnosis of why
it cannot be pushed materially further on the existing sensor channels:
two anomalous training sets account for most of the LOSO error, and
the distribution shift to RVS\,304 caps val performance. This is enough
to support the project's main argument --- that physics-informed loss
(Stage 3) is a better lever than additional sensor work --- and to
provide a clean reference point for the sensor-fusion comparisons in
the multimodal sub-track.

\end{document}
